\section{The families, drawn}
\label{app:catalog}

Table~\ref{tab:catalog} gives the fourteen families as algebra.
Figure~\ref{fig:plate} gives them as pictures, each one alone and then superposed
with a copy of itself under a small perturbation, so that the fringe system each
one generates can be seen next to the family that generates it.

The perturbation differs by family, and which one is natural is itself
informative. Families with an orientation (the two line families, the three
lattices, the three open curves) fringe under a small \emph{rotation}, a few
degrees, because rotating them changes the direction of $\nabla\idx$ without
changing its magnitude. Closed families (the four concentric shapes and the
spiral) fringe under a small \emph{detuning} of the pitch instead, five or six
percent, because rotating a concentric family about its own center does nothing at
all. This is not a fact about our implementation. It is $\het$ of
Eq.~\ref{eq:ratio} deciding what counts as a perturbation: for a family whose
index gradient is radial, rotation about the center leaves $\nabla\idx$ unchanged
and so leaves $\het=0$, and there is nothing to see.

The log spiral takes this to its logical end, and the carrier classification
of \S\ref{sec:taxonomy} says it must: the family is a group orbit of the
loxodromic flow, and a group family is immune to its own flow. Rotating a log
spiral is a phase shift; rescaling it --- the detune that fringes every other
closed family --- is another, because scaling adds a constant to $\ln r$. A
similarity can do nothing visible to it. The one perturbation left to a plate
of similarities is the one thing a similarity cannot reach, the handedness,
so its cell shows the mirror pair: a self-similar rosette of lenses,
$2M$ around, telescoping into the center without a characteristic scale.

\newlength{\platew}
% Three rows of cells plus a full-width caption have to clear \textheight, or
% the float is never placed and silently vanishes (it did: at 0.178 with the
% caption squeezed into the last cell slot, row three outgrew the page). A
% cell is the panel plus a hairline, 2pt and one \footnotesize line; at 0.168
% that is ~188pt, so 3 rows, two 7pt gaps, and a four-line caption come in
% just under the 615pt available.
\setlength{\platew}{0.168\textwidth}
% One cell of the catalog plate: a family alone above, beating below, and a label
% naming the perturbation that produced the beat.
\newcommand{\platecell}[3]{%
  \begin{minipage}[t]{\platew}%
    \vspace{0pt}% align cell tops, not first baselines
    \centering
    \setlength{\fboxsep}{0pt}\setlength{\fboxrule}{0.3pt}%
    \fcolorbox{soft!45}{white}{%
      \includegraphics[width=\dimexpr\platew-2\fboxrule\relax]{figures/plate-#2.png}}%
    \\[2pt]
    {\footnotesize #1\,\ #3}%
  \end{minipage}%
}
\newcommand{\platedeg}[1]{\textcolor{soft}{$\cdot$\,$#1^{\circ}$}}
\newcommand{\platedet}[1]{\textcolor{soft}{$\cdot\,\times#1$}}
\newcommand{\platemir}{\textcolor{soft}{$\cdot$\,mirror}}
% A row is one unbreakable hbox. Left to a paragraph, a row that is a hair too wide
% wraps its last cell and the float grows by a whole row.
\newcommand{\platerow}[1]{\par\noindent\hbox to \linewidth{#1}}

\begin{figure*}[p]
  \Description{Fourteen cells in a grid, each a pair of stacked black-on-white line
  drawings. The upper drawing of a cell is one curve family: parallel lines, a
  radial pencil with a blank disc at its hub, concentric circles, squares, triangles,
  hexagons, a square lattice, a hexagonal lattice, a triangular lattice, a wave
  train, a nest of parabolae, a nest of hyperbolae, a spiral, a logarithmic spiral.
  The lower drawing
  superposes that family with a slightly rotated or slightly rescaled copy of itself,
  and in every case a second, coarser pattern of bands appears that is not present
  above.}
  \centering
  \platerow{%
    \platecell{1}{parallel}{parallel \platedeg{4}}\hfil
    \platecell{2}{radial}{radial \platedeg{5}}\hfil
    \platecell{3}{circles}{circles \platedet{1.06}}\hfil
    \platecell{4}{squares}{squares \platedet{1.06}}\hfil
    \platecell{5}{triangles}{triangles \platedet{1.06}}}

  \vskip 7pt
  \platerow{%
    \platecell{6}{hexagons}{hexagons \platedet{1.06}}\hfil
    \platecell{7}{grid-square}{square grid \platedeg{5}}\hfil
    \platecell{8}{grid-hex}{hex grid \platedeg{5}}\hfil
    \platecell{9}{grid-triangle}{tri.\ grid \platedeg{5}}\hfil
    \platecell{10}{wave}{wave \platedeg{3}}}

  \vskip 7pt
  \platerow{%
    \platecell{11}{parabola}{parabola \platedeg{3}}\hfil
    \platecell{12}{hyperbola}{hyperbola \platedeg{4}}\hfil
    \platecell{13}{spiral}{spiral \platedet{1.05}}\hfil
    \platecell{14}{log}{log spiral \platemir}\hfil
    \hspace{\platew}}

  \vskip 4pt
  \caption{\figlead{Fourteen families, alone and beating.} Each cell is one
  family: above, by itself from its distance field; below, superposed with a
  perturbed copy, so the fringe system sits directly under the family that
  generates it. The label names the perturbation --- degrees of rotation for
  oriented families, a pitch multiplier for the closed ones, a mirror for the
  log spiral --- which is $\het$ of Eq.~\ref{eq:ratio} deciding what counts as
  a perturbation, not a choice of ours (\S\ref{app:catalog}). Every panel
  spans $250$ world units at one stroke width; none is a crop or resample of
  another.}
  \label{fig:plate}
\end{figure*}

Three details in the plate repay a second look. The radial pencil's core is
genuinely dense, and \emph{Start} cuts a disc out of it rather than fading it.
The hexagonal grid's edges are hexagon sides, computed as a support form over
six normals, not three superposed line families. And the hyperbola family
carries no member through the origin, because its $n=0$ level set is the
degenerate asymptote pair.

Two families deserve a note on their parameterization, because the natural
parameter and the useful parameter differ.

\paragraph{The parabola's bend.} The family is $y=ax^2+ns+\varphi$, and $a$ is
exposed to the author as a slider $B$ with $a=B/100$. The scaling is not
cosmetic. Perpendicular member spacing is $s/\lVert\nabla\ph\rVert =
s/\sqrt{1+4a^2x^2}$, which falls off linearly in $\lvert x\rvert$, so a value of
$a$ that looks reasonable at the origin has closed the family to nothing by the
edge of the frame. On a $250$-unit panel, $a=0.5$, an unremarkable-looking
number, packs several hundred members into the corners. The slider therefore
spans a range that keeps the whole frame resolvable, and the useful settings all
have $B<1$.

\paragraph{The spiral's pitch.} Archimedean spirals are $r=b\theta$, and a
family of them at radial interval $s$ has $M=\mathrm{round}(\lvert B\rvert/s)$
arms for a pitch slider $B$. The rounding is load-bearing. $\ph=r-b\theta$ has a
branch cut wherever $\mathrm{atan2}$ does, and the level sets close across it only
if $b\,2\pi$ is an integer multiple of $s$, that is, only if $M$ is an integer.
Pass $B$ through unrounded and every member of the family has a visible seam along
the negative $x$ axis. Setting $B=0$ gives $M=1$ but $b=0$, which is concentric
circles, so the slider degrades to the family it should degrade to. The log
spiral inherits all of it verbatim --- the rounding, the seam argument, and the
handedness in the sign --- with $r$ replaced by $32\ln r$, and its $B=0$
degrades to geometrically spaced rings the same way.
